\section{Revisão crítica da formulação do problema}
Diferenciem sintomas, possíveis causas e soluções já imaginadas. Causas sem evidência devem permanecer identificadas como hipóteses.

\begin{table}[!ht]\centering\small\renewcommand{\arraystretch}{1.35}
\begin{tabularx}{\linewidth}{|>{\raggedright\arraybackslash}X|>{\raggedright\arraybackslash}X|>{\raggedright\arraybackslash}X|}
\hline
\textbf{Sintomas / manifestações} &
\textbf{Possíveis causas / hipóteses} &
\textbf{Soluções já imaginadas} \\
\hline
Avisos não chegam aos destinatários ou chegam com atraso &
Utilização de múltiplos canais de comunicação sem coordenação central &
Centralizar a divulgação das informações em um único ambiente \\
\hline
Mensagens importantes se perdem em grupos de WhatsApp &
Excesso de mensagens e ausência de categorização adequada &
Organização dos comunicados por tipo e público-alvo \\
\hline
Pais possuem dificuldades para acompanhar a rotina escolar dos filhos &
Dependência de canais informais e descentralizados &
Disponibilizar acompanhamento integrado de informações escolares \\
\hline
Direção e coordenação têm dificuldade em controlar a divulgação dos comunicados &
Ausência de fluxo formal de publicação e aprovação &
Estabelecer processo de moderação e validação dos avisos \\
\hline
Não há confirmação sobre quem recebeu ou visualizou informações importantes &
Falta de mecanismos de acompanhamento da comunicação &
Registro de recebimento e leitura das mensagens \\
\hline
\end{tabularx}
\end{table}

\begin{notebox}[Teste sem tecnologia]
Reescrevam o problema sem mencionar aplicativo, plataforma, sistema, IA, banco de dados ou outra tecnologia.
\end{notebox}

A escola possui dificuldades para organizar e distribuir informações de forma eficiente entre direção, coordenação, professores, alunos e responsáveis. Os canais atualmente utilizados não garantem que os comunicados alcancem os destinatários corretos no momento adequado, comprometendo o acompanhamento das atividades escolares, a participação das famílias e a efetividade da comunicação institucional.

\begin{notebox}[Verificação]
Se retirarmos a solução tecnológica, o problema continua existindo?

\noindent $\boxtimes$ Sim \qquad $\square$ Não --- reformular.
\end{notebox}
